\documentclass[12pt,a4paper]{article}

% 中文支持
\usepackage[UTF8]{ctex}
\usepackage{xeCJK}

% 页面设置
\usepackage[margin=2.5cm]{geometry}
\usepackage{fancyhdr}
\pagestyle{fancy}
\fancyhf{}
\fancyhead[C]{图片管理系统实验报告}
\fancyfoot[C]{\thepage}

% 数学和符号
\usepackage{amsmath,amsfonts,amssymb}
\usepackage{bm}

% 图形和颜色
\usepackage{graphicx}
\usepackage{xcolor}
\usepackage{tikz}
\usetikzlibrary{shapes,arrows,positioning,decorations.pathreplacing}

% 代码高亮
\usepackage{listings}
\usepackage{xcolor}

% 漂亮的框
\usepackage{tcolorbox}
\tcbuselibrary{listings,breakable,skins}

% 表格
\usepackage{booktabs}
\usepackage{tabularx}
\usepackage{array}

% 超链接
\usepackage{hyperref}
\hypersetup{
    colorlinks=true,
    linkcolor=blue,
    filecolor=magenta,      
    urlcolor=cyan,
    pdftitle={图片管理系统实验报告},
    pdfauthor={实验者},
}

% 代码样式
\definecolor{codegreen}{rgb}{0,0.6,0}
\definecolor{codegray}{rgb}{0.5,0.5,0.5}
\definecolor{codepurple}{rgb}{0.58,0,0.82}
\definecolor{backcolour}{rgb}{0.95,0.95,0.92}

\lstdefinestyle{mystyle}{
    backgroundcolor=\color{backcolour},   
    commentstyle=\color{codegreen},
    keywordstyle=\color{magenta},
    numberstyle=\tiny\color{codegray},
    stringstyle=\color{codepurple},
    basicstyle=\ttfamily\footnotesize,
    breakatwhitespace=false,         
    breaklines=true,                 
    captionpos=b,                    
    keepspaces=true,                 
    numbers=left,                    
    numbersep=5pt,                  
    showspaces=false,                
    showstringspaces=false,
    showtabs=false,                  
    tabsize=2,
    language=Python
}

\lstset{style=mystyle}

% 自定义命令
\newcommand{\tech}[1]{\textbf{\textcolor{blue}{#1}}}
\newcommand{\feature}[1]{\textcolor{green!60!black}{✓ #1}}

% 标题信息
\title{\textbf{\Huge 图片管理系统\\实验报告}}
\author{
    \Large 实验者姓名 \\
    \normalsize 学号：XXXXXXXXXX \\
    \normalsize 日期：\today
}
\date{}

\begin{document}

\maketitle
\thispagestyle{empty}

\newpage
\tableofcontents
\newpage

\section{项目概述}

\subsection{项目背景}
随着数字图像数量的快速增长，传统的文件系统管理方式已无法满足用户对图片组织、检索和管理的需求。本项目旨在开发一个功能完善的图片管理系统，集成人工智能技术，实现智能化的图片管理和分析。

\subsection{项目目标}
\begin{itemize}
    \item 实现图片的上传、存储、浏览和管理功能
    \item 集成AI视觉模型，实现自动标签生成和图像描述
    \item 提供友好的用户界面和良好的用户体验
    \item 支持多用户权限管理和数据安全
    \item 实现图片编辑、批量操作等高级功能
\end{itemize}

\section{技术栈}

\subsection{后端技术栈}

\begin{table}[h]
\centering
\caption{后端核心技术栈}
\begin{tabularx}{\textwidth}{|X|X|X|}
\hline
\textbf{技术/框架} & \textbf{版本} & \textbf{用途} \\
\hline
Python & 3.11 & 主要开发语言 \\
\hline
Flask & 最新版 & Web框架，提供RESTful API \\
\hline
Connexion & 最新版 & OpenAPI规范支持，自动生成API文档 \\
\hline
SQLAlchemy & 最新版 & ORM框架，数据库操作抽象层 \\
\hline
Flask-SQLAlchemy & 最新版 & Flask的SQLAlchemy集成 \\
\hline
PyMySQL & 最新版 & MySQL数据库驱动 \\
\hline
Pillow (PIL) & 最新版 & 图像处理库 \\
\hline
OpenCV & 最新版 & 计算机视觉库，图像处理 \\
\hline
PyYAML & 最新版 & 配置文件解析 \\
\hline
bcrypt & 最新版 & 密码加密 \\
\hline
PyJWT & 最新版 & JWT令牌生成和验证 \\
\hline
requests & 最新版 & HTTP客户端，调用外部API \\
\hline
exifread & 最新版 & EXIF数据提取 \\
\hline
\end{tabularx}
\end{table}

\subsection{前端技术栈}

\begin{table}[h]
\centering
\caption{前端核心技术栈}
\begin{tabularx}{\textwidth}{|X|X|X|}
\hline
\textbf{技术/框架} & \textbf{版本} & \textbf{用途} \\
\hline
React & 18.2.0 & 前端UI框架 \\
\hline
TypeScript & 5.2.2 & 类型安全的JavaScript超集 \\
\hline
Vite & 4.4.11 & 现代化构建工具 \\
\hline
Material-UI (MUI) & 5.14.16 & React组件库 \\
\hline
React Router & 6.16.0 & 前端路由管理 \\
\hline
Axios & 1.5.1 & HTTP客户端 \\
\hline
React Dropzone & 14.2.3 & 文件拖拽上传组件 \\
\hline
React Infinite Scroll & 6.1.0 & 无限滚动加载 \\
\hline
Notistack & 3.0.1 & 通知消息管理 \\
\hline
\end{tabularx}
\end{table}

\subsection{数据库技术}

\begin{itemize}
    \item \tech{MySQL 8.0}：关系型数据库，存储用户、图片、标签等结构化数据
    \item \tech{SQLite}：可选数据库，用于开发测试环境
\end{itemize}

\subsection{AI技术}

\begin{itemize}
    \item \tech{Qwen3 VL (通义千问视觉模型)}：阿里云DashScope提供的多模态大语言模型
    \item \tech{模型版本}：qwen3-vl-flash（快速响应版本）
    \item \tech{API接口}：OpenAI兼容模式的RESTful API
    \item \tech{功能}：
    \begin{itemize}
        \item 自动图像标签生成（最多5个标签）
        \item 图像内容描述生成（50-200字）
    \end{itemize}
\end{itemize}

\subsection{容器化技术}

\begin{itemize}
    \item \tech{Docker}：应用容器化
    \item \tech{Docker Compose}：多容器应用编排
    \item 服务包括：
    \begin{itemize}
        \item 后端服务容器（Python Flask）
        \item 前端服务容器（Node.js + Vite）
        \item 数据库服务容器（MySQL 8.0）
    \end{itemize}
\end{itemize}

\section{系统架构}

\subsection{整体架构}

\begin{figure}[h]
\centering
\begin{tikzpicture}[
    node distance=1.5cm,
    auto,
    thick,
    main node/.style={rectangle, draw, fill=blue!20, text width=3cm, text centered, rounded corners, minimum height=1cm},
    db node/.style={cylinder, draw, fill=green!20, text width=2.5cm, text centered, minimum height=1cm},
    api node/.style={rectangle, draw, fill=orange!20, text width=2.5cm, text centered, rounded corners, minimum height=1cm}
]

% 前端层
\node[main node] (frontend) {前端 (React)};
\node[main node, right=of frontend] (nginx) {Vite Dev Server};

% 后端层
\node[main node, below=of frontend] (backend) {后端 (Flask)};
\node[api node, right=of backend] (qwen) {Qwen3 VL API};

% 数据库层
\node[db node, below=of backend] (mysql) {MySQL 8.0};

% 存储层
\node[main node, below=of mysql] (storage) {文件存储系统};

% 连接
\draw[->] (frontend) -- (nginx);
\draw[->] (nginx) -- node {HTTP/API} (backend);
\draw[->] (backend) -- node {SQL} (mysql);
\draw[->] (backend) -- node {REST API} (qwen);
\draw[->] (backend) -- node {文件I/O} (storage);

\end{tikzpicture}
\caption{系统整体架构图}
\end{figure}

\subsection{数据流图}

\begin{figure}[h]
\centering
\begin{tikzpicture}[
    node distance=2cm,
    auto,
    thick,
    process/.style={rectangle, draw, fill=cyan!20, text width=2.5cm, text centered, minimum height=1cm},
    data/.style={ellipse, draw, fill=yellow!20, text width=2cm, text centered}
]

\node[process] (upload) {图片上传};
\node[process, right=of upload] (process) {图像处理};
\node[process, right=of process] (ai) {AI分析};
\node[process, below=of process] (db) {数据库存储};
\node[data, below=of upload] (file) {文件系统};

\draw[->] (upload) -- (process);
\draw[->] (process) -- (ai);
\draw[->] (process) -- (db);
\draw[->] (upload) -- (file);
\draw[->] (ai) -- (db);

\end{tikzpicture}
\caption{图片处理数据流图}
\end{figure}

\section{核心功能实现}

\subsection{用户认证与权限管理}

\subsubsection{实现技术}
\begin{itemize}
    \item \tech{JWT (JSON Web Token)}：无状态身份认证
    \item \tech{bcrypt}：密码哈希加密
    \item \tech{角色权限系统}：pending、user、admin三级权限
\end{itemize}

\subsubsection{功能特性}
\begin{itemize}
    \item \feature{用户注册}：邮箱验证，初始状态为pending
    \item \feature{用户登录}：JWT令牌生成，24小时有效期
    \item \feature{权限管理}：管理员可审批用户、修改角色
    \item \feature{默认管理员}：系统初始化时自动创建admin账户
\end{itemize}

\subsection{图片管理功能}

\subsubsection{图片上传}
\begin{itemize}
    \item \feature{多文件上传}：支持同时上传多张图片
    \item \feature{文件验证}：格式验证（JPG、PNG、GIF、BMP、WebP）
    \item \feature{安全存储}：随机哈希文件名，防止冲突
    \item \feature{EXIF提取}：自动提取拍摄时间、GPS位置等信息
    \item \feature{缩略图生成}：自动生成300px缩略图，Base64编码存储
\end{itemize}

\subsubsection{图片浏览}
\begin{itemize}
    \item \feature{无限滚动}：React Infinite Scroll组件实现懒加载
    \item \feature{分页查询}：后端分页，每页20张图片
    \item \feature{标签筛选}：按标签筛选图片
    \item \feature{搜索功能}：按文件名搜索
\end{itemize}

\subsubsection{图片编辑}
\begin{itemize}
    \item \feature{图片裁剪}：支持四边裁剪（百分比）
    \item \feature{色彩调整}：亮度、对比度调整
    \item \feature{编辑历史}：自动备份原始文件
    \item \feature{缩略图更新}：编辑后自动更新缩略图
\end{itemize}

\subsubsection{批量操作}
\begin{itemize}
    \item \feature{批量选择}：复选框多选，支持全选
    \item \feature{批量删除}：一次性删除多张图片
    \item \feature{软删除}：标记删除，可恢复
    \item \feature{权限检查}：只能删除自己的图片（管理员除外）
\end{itemize}

\subsection{标签管理系统}

\subsubsection{功能特性}
\begin{itemize}
    \item \feature{标签创建}：手动创建自定义标签
    \item \feature{标签关联}：多对多关系，一张图片可有多个标签
    \item \feature{标签管理}：添加/移除图片标签
    \item \feature{标签筛选}：按标签查看相关图片
\end{itemize}

\subsection{AI智能分析功能}

\subsubsection{Qwen3 VL集成}
\begin{itemize}
    \item \feature{自动标签生成}：上传时自动调用API生成标签
    \item \feature{图像描述生成}：生成50-200字的图像描述
    \item \feature{错误处理}：API调用失败不影响图片上传
    \item \feature{重试机制}：最多重试3次，指数退避
    \item \feature{配置管理}：通过config.yaml配置API密钥和参数
\end{itemize}

\subsubsection{技术实现细节}
\begin{tcolorbox}[colback=blue!5!white,colframe=blue!75!black,title=AI标签生成流程]
\begin{enumerate}
    \item 图片上传后，转换为Base64编码
    \item 调用Qwen3 VL API，发送图片和提示词
    \item 解析API响应，提取标签文本
    \item 标签文本分割（支持中文逗号、英文逗号等）
    \item 验证标签有效性（长度2-20字符）
    \item 保存标签到数据库，关联到图片
\end{enumerate}
\end{tcolorbox}

\section{数据库设计}

\subsection{数据模型}

\begin{table}[h]
\centering
\caption{核心数据表结构}
\begin{tabularx}{\textwidth}{|l|X|}
\hline
\textbf{表名} & \textbf{主要字段} \\
\hline
users & id, username, email, password\_hash, role, created\_at \\
\hline
images & id, filename, file\_path, file\_hash, file\_size, width, height, \\
 & taken\_at, location\_lat, location\_lng, thumbnail, description, \\
 & is\_deleted, user\_id, created\_at \\
\hline
tags & id, name, created\_at \\
\hline
image\_tags & image\_id, tag\_id (多对多关系表) \\
\hline
image\_edits & id, image\_id, operation, backup\_path, edit\_params, created\_at \\
\hline
\end{tabularx}
\end{table}

\subsection{关系设计}
\begin{itemize}
    \item \textbf{用户-图片}：一对多关系（一个用户可有多张图片）
    \item \textbf{图片-标签}：多对多关系（一张图片可有多个标签）
    \item \textbf{图片-编辑历史}：一对多关系（一张图片可有多次编辑记录）
\end{itemize}

\section{API设计}

\subsection{RESTful API端点}

\begin{table}[h]
\centering
\caption{主要API端点}
\begin{tabularx}{\textwidth}{|l|l|X|}
\hline
\textbf{方法} & \textbf{路径} & \textbf{功能} \\
\hline
POST & /api/v1/auth/login & 用户登录 \\
POST & /api/v1/auth/register & 用户注册 \\
GET & /api/v1/users & 获取用户列表（管理员） \\
POST & /api/v1/users/<id>/approve & 审批用户 \\
PUT & /api/v1/users/<id>/role & 修改用户角色 \\
POST & /api/v1/images/upload & 上传图片 \\
GET & /api/v1/images & 获取图片列表 \\
GET & /api/v1/images/<id> & 获取图片详情 \\
POST & /api/v1/images/<id>/edit & 编辑图片 \\
POST & /api/v1/images/<id>/delete & 删除图片 \\
POST & /api/v1/images/batch-delete & 批量删除图片 \\
POST & /api/v1/images/<id>/restore & 恢复图片 \\
GET & /api/v1/tags & 获取标签列表 \\
POST & /api/v1/tags & 创建标签 \\
POST & /api/v1/images/<id>/tags & 添加标签到图片 \\
DELETE & /api/v1/images/<id>/tags & 从图片移除标签 \\
GET & /api/v1/config & 获取配置 \\
PUT & /api/v1/config & 更新配置 \\
POST & /api/v1/config/test-qwen & 测试Qwen连接 \\
\hline
\end{tabularx}
\end{table}

\section{关键技术实现}

\subsection{图像处理技术}

\subsubsection{缩略图生成}
使用Pillow库实现高效的缩略图生成：
\begin{lstlisting}[language=Python]
def generate_thumbnail(image_path, max_size=300):
    with Image.open(image_path) as img:
        img.thumbnail((max_size, max_size), Image.Resampling.LANCZOS)
        buffer = io.BytesIO()
        img.save(buffer, format='JPEG', quality=85)
        return base64.b64encode(buffer.getvalue()).decode('utf-8')
\end{lstlisting}

\subsubsection{EXIF数据提取}
自动提取图片的元数据信息：
\begin{itemize}
    \item 拍摄时间（taken\_at）
    \item GPS坐标（location\_lat, location\_lng）
    \item 相机信息、拍摄参数等
\end{itemize}

\subsection{AI API集成}

\subsubsection{Qwen3 VL客户端实现}
\begin{tcolorbox}[colback=green!5!white,colframe=green!75!black,title=关键特性]
\begin{itemize}
    \item 支持OpenAI兼容模式和标准DashScope模式
    \item 自动图片压缩（最大1024x1024）
    \item Base64编码传输
    \item 智能响应解析（兼容多种响应格式）
    \item 完善的错误处理和重试机制
    \item 超时控制（默认30秒）
\end{itemize}
\end{tcolorbox}

\subsection{前端技术亮点}

\subsubsection{无限滚动实现}
使用React Infinite Scroll Component实现流畅的图片加载体验：
\begin{itemize}
    \item 自动检测滚动位置
    \item 懒加载下一页数据
    \item 加载状态提示
    \item 性能优化（虚拟滚动）
\end{itemize}

\subsubsection{拖拽上传}
使用React Dropzone实现友好的文件上传体验：
\begin{itemize}
    \item 支持拖拽文件到指定区域
    \item 多文件选择
    \item 文件类型验证
    \item 上传进度显示
\end{itemize}

\section{系统部署}

\subsection{Docker容器化}

\subsubsection{容器配置}
\begin{itemize}
    \item \textbf{后端容器}：Python 3.11-slim基础镜像
    \item \textbf{前端容器}：Node.js + Vite开发服务器
    \item \textbf{数据库容器}：MySQL 8.0官方镜像
\end{itemize}

\subsubsection{Docker Compose配置}
\begin{lstlisting}[language=yaml]
services:
  db:
    image: mysql:8.0
    environment:
      MYSQL_ROOT_PASSWORD: rootpassword
      MYSQL_DATABASE: image_manager
  
  backend:
    build:
      dockerfile: Dockerfile.backend
    ports:
      - "8080:8080"
    depends_on:
      - db
  
  frontend:
    build:
      dockerfile: Dockerfile.frontend
    ports:
      - "3000:3000"
\end{lstlisting}

\section{测试与验证}

\subsection{功能测试}

\begin{table}[h]
\centering
\caption{功能测试结果}
\begin{tabularx}{\textwidth}{|X|c|}
\hline
\textbf{测试项目} & \textbf{结果} \\
\hline
用户注册和登录 & ✓ 通过 \\
图片上传（单文件/多文件） & ✓ 通过 \\
图片浏览和分页 & ✓ 通过 \\
图片编辑（裁剪/色彩调整） & ✓ 通过 \\
批量删除功能 & ✓ 通过 \\
标签管理 & ✓ 通过 \\
AI标签生成 & ✓ 通过 \\
AI描述生成 & ✓ 通过 \\
权限控制 & ✓ 通过 \\
API接口测试 & ✓ 通过 \\
\hline
\end{tabularx}
\end{table}

\subsection{性能测试}

\begin{itemize}
    \item \textbf{图片上传}：单张图片（2MB）上传时间 < 3秒（含AI处理）
    \item \textbf{图片列表加载}：20张图片加载时间 < 1秒
    \item \textbf{AI标签生成}：平均响应时间 2-5秒
    \item \textbf{并发处理}：支持多用户同时上传
\end{itemize}

\section{项目亮点}

\subsection{技术创新}
\begin{enumerate}
    \item \textbf{AI深度集成}：将Qwen3 VL视觉模型无缝集成到图片管理流程中
    \item \textbf{智能标签系统}：自动生成标签，减少用户手动标注工作量
    \item \textbf{图像描述生成}：自动生成图像描述，提升图片检索能力
    \item \textbf{完善的错误处理}：AI功能失败不影响核心功能使用
\end{enumerate}

\subsection{用户体验}
\begin{enumerate}
    \item \textbf{现代化UI}：Material-UI设计，界面美观易用
    \item \textbf{响应式设计}：适配不同屏幕尺寸
    \item \textbf{实时反馈}：上传进度、操作提示等
    \item \textbf{批量操作}：提高管理效率
\end{enumerate}

\subsection{代码质量}
\begin{enumerate}
    \item \textbf{模块化设计}：前后端分离，职责清晰
    \item \textbf{类型安全}：TypeScript + Python类型注解
    \item \textbf{错误处理}：完善的异常捕获和日志记录
    \item \textbf{代码规范}：遵循PEP 8和ESLint规范
\end{enumerate}

\section{总结与展望}

\subsection{项目总结}

本项目成功实现了一个功能完善的图片管理系统，集成了AI视觉分析能力，提供了良好的用户体验。系统采用现代化的技术栈，实现了前后端分离、容器化部署，具有良好的可扩展性和可维护性。

\subsection{技术收获}

\begin{itemize}
    \item 掌握了Flask RESTful API设计和开发
    \item 学习了React + TypeScript前端开发
    \item 了解了Docker容器化部署
    \item 掌握了AI API集成和调用
    \item 学习了数据库设计和ORM使用
    \item 提升了全栈开发能力
\end{itemize}

\subsection{未来改进方向}

\begin{enumerate}
    \item \textbf{性能优化}：
    \begin{itemize}
        \item 图片CDN加速
        \item 数据库查询优化
        \item 前端代码分割和懒加载
    \end{itemize}
    
    \item \textbf{功能扩展}：
    \begin{itemize}
        \item 图片相似度搜索
        \item 人脸识别和分类
        \item 图片自动分类（按场景、对象等）
        \item 移动端APP开发
    \end{itemize}
    
    \item \textbf{安全性增强}：
    \begin{itemize}
        \item HTTPS支持
        \item 文件病毒扫描
        \item 更严格的权限控制
        \item 数据加密存储
    \end{itemize}
\end{enumerate}

\section{参考文献}

\begin{enumerate}
    \item Flask官方文档. \url{https://flask.palletsprojects.com/}
    \item React官方文档. \url{https://react.dev/}
    \item Qwen3 VL模型文档. \url{https://help.aliyun.com/zh/model-studio/}
    \item Docker官方文档. \url{https://docs.docker.com/}
    \item Material-UI文档. \url{https://mui.com/}
\end{enumerate}

\newpage
\section*{附录}

\subsection*{A. 项目文件结构}

\begin{lstlisting}[language=bash, basicstyle=\ttfamily\small]
image-manage/
├── backend/              # 后端代码
│   ├── app.py           # Flask应用入口
│   ├── database.py      # 数据库模型
│   ├── qwen_client.py   # Qwen API客户端
│   ├── controllers/     # 控制器
│   └── utils.py         # 工具函数
├── frontend/            # 前端代码
│   ├── src/
│   │   ├── pages/       # 页面组件
│   │   ├── services/    # API服务
│   │   └── contexts/     # React上下文
│   └── package.json     # 依赖配置
├── scripts/             # 脚本文件
│   └── run_backend.py   # 后端启动脚本
├── docker-compose.yml   # Docker编排配置
├── Dockerfile.backend   # 后端镜像配置
├── Dockerfile.frontend  # 前端镜像配置
├── requirements.txt     # Python依赖
└── config.yaml          # 系统配置
\end{lstlisting}

\subsection*{B. 核心代码示例}

\textbf{B.1 图片上传处理}

\begin{lstlisting}[language=Python]
def upload_images():
    files = request.files.getlist('files')
    for file in files:
        # 保存文件
        storage_path = get_storage_path(base_path, filename, file_hash)
        save_uploaded_file(file, storage_path)
        
        # 提取EXIF信息
        exif_data = extract_exif_data(storage_path)
        
        # 生成缩略图
        thumbnail = generate_thumbnail(storage_path)
        
        # AI标签生成
        qwen_client = QwenVLClient(config)
        ai_tags = qwen_client.analyze_image(storage_path)
        
        # AI描述生成
        ai_description = qwen_client.generate_description(storage_path)
        
        # 保存到数据库
        image = Image(...)
        session.add(image)
\end{lstlisting}

\textbf{B.2 批量删除实现}

\begin{lstlisting}[language=Python]
def batch_delete_images():
    image_ids = request.json.get('image_ids', [])
    images = session.query(Image).filter(
        Image.id.in_(image_ids),
        Image.is_deleted == False
    ).all()
    
    for image in images:
        image.is_deleted = True
        # 记录删除操作
        edit_record = ImageEdit(...)
        session.add(edit_record)
    
    session.commit()
\end{lstlisting}

\end{document}

